% !PRINTED 52
\weberpage{52}
%<<CONT
gemäss bestimmen. Ist $\alpha_1, \alpha_2$ eine Lösung dieser
Congruenz, in der nicht beide Zahlen Null sind, so erhalten wir
alle Lösungen, wenn wir in $h\alpha_1, h\alpha_2$ den Factor $h$
die Reihe der Zahlen $0, 1 \dots p-1$ durchlaufen lassen, und die
Gruppe $T$ wird also, wenn $\alpha_1, \alpha_2$ ein festes, der
Bedingung (9) genügendes Werthpaar ist,
\[
(A_1^{\alpha_1}A_2^{\alpha_2})^h, \quad h = 0, 1, \dots p-1.
\]

Die reciproke Gruppe $U$ erhält man, wenn man alle Werthe der
$\beta$ sucht, die der Bedingung (9) genügen, und man findet also
die Gruppe $U$ in der Form
\[
(A_1^{\beta_1}A_2^{\beta_2})^h, \quad h = 0, 1, \dots p-1,
\]
worin $\alpha_1, \alpha_2, \beta_1, \beta_2$ jetzt vier feste, der
Bedingung (9) genügende Zahlen sind. Setzt man $\beta_1 = 0$,
$\beta_2 = 1$, $\alpha_1 = 1$, $\alpha_2 = 0$, so erhält man den
besonderen Fall der beiden reciproken Gruppen
\[
A_1^h, \quad A_2^h, \quad h = 0, 1, \dots p-1.
\]

\begin{center}
\S.\,13.

\gesperrt{Die Geschlechter in einer Abel'schen Gruppe.}
\end{center}

Ein Element einer \gesperrt{Abel}'schen Gruppe $S$, das mit seinem
entgegengesetzten identisch ist, dessen zweite Potenz also gleich
dem Einheitselemente ist, wird ein \gesperrt{zweiseitiges
Element}\footnote{Vergl. Bd.~I, S.~390, Anmerkung.} genannt. Das
Einheitselement gehört also immer zu den zweiseitigen.

Stellen wir die Elemente von $S$ durch eine Basis
$A_1, A_2 \dots A_\nu$ dar, deren Elemente die Grade
$a_1, a_2 \dots a_\nu$ haben, so wird
\weq{1}{A = A_1^{\alpha_1}A_2^{\alpha_2} \dots A_\nu^{\alpha_\nu}}
dann und nur dann ein zweiseitiges Element sein, wenn
\weq{2}{2\alpha_1 \equiv 0\ (\mathrm{mod}\ a_1), \ 2\alpha_2 \equiv 0\ (\mathrm{mod}\ a_2), \dots 2\alpha_\nu \equiv 0\ (\mathrm{mod}\ a_\nu).}
Wenn nun unter den Invarianten der Gruppe $\lambda$ mal eine
Potenz von 2 vorkommt, so können wir die Basis von $S$ so
geordnet annehmen, dass $a_1, a_2 \dots a_\lambda$ gerade,
$a_{\lambda+1} \dots a_\nu$ ungerade Zahlen sind. Dann ergeben
sich die Lösungen von (2):
%CONT>>
